\chapter{Template Overview}
\label{ch:overview}

As noted in the introduction, this document class can be used to prepare
many different renderings of the same thesis --- a copy for review, a
clean copy for printing and binding, and a copy for reading on screen ---
all by selecting the appropriate {\itshape template style} and {\itshape
  template parameters}.

This document explains the major features of the document class. For
further information on the original ACM class, the {\itshape \LaTeX\
User's Guide} is available from
\url{https://www.acm.org/publications/proceedings-template}.

\section{Template Styles}
\label{sec:styles}

The primary parameter given to the document class is the {\itshape
  template style}, which corresponds to the intended use of the output.
This parameter is enclosed in square brackets and is a part of the
\verb|documentclass| command:
\begin{verbatim}
  \documentclass[STYLE]{thesis}
\end{verbatim}

The thesis template defines three styles:
\begin{itemize}
\item {\texttt{print}}: a two-sided (duplex) copy with binding margins,
  suitable for printing and binding. This is the default.
\item {\texttt{screen}}: a one-sided copy for reading on screen, with
  coloured hyperlinks.
\item {\texttt{review}}: a one-sided reading copy with line numbers.
\end{itemize}

\section{Template Parameters}
\label{sec:parameters}

In addition to specifying the {\itshape template style}, there are a
number of {\itshape template parameters} which modify some part of the
applied template style. The base font size is 14\,pt by default, and
10, 11, 12, 17 and 20\,pt are also accepted. Frequently-used parameters
include:
\begin{itemize}
\item {\texttt{onehalf}}, {\texttt{double}}: line spacing.
\item {\texttt{font=times}}: a Times-like text face ({\texttt{font=libertine}}
  is the default).
\item {\texttt{authoryear}}: author--year citations.
\item {\texttt{biblatex}}: use \BibLaTeX\ with the \texttt{acmnumeric}
  or \texttt{acmauthoryear} styles instead of \BibTeX.
\end{itemize}


Testing the listing
\begin{itemize}
  \item Something
  \item Another one
\end{itemize}

This document uses the following string as the first command in the
source file:
\begin{verbatim}
\documentclass[print]{thesis}
\end{verbatim}

\section{Files and Building}
\label{sec:files}

The template is built with the usual \LaTeX\ tools. With \texttt{latexmk}:
\begin{verbatim}
  latexmk -pdf thesis.tex
\end{verbatim}
or, manually, \verb|pdflatex|, then \verb|bibtex|, then \verb|pdflatex|
twice more. A full \TeX\ Live (or MiKTeX) installation is required; the
class uses the standard packages \texttt{extsizes}, \texttt{libertine},
\texttt{newtxmath}, \texttt{geometry}, \texttt{titlesec},
\texttt{fancyhdr}, \texttt{setspace}, \texttt{natbib}, \texttt{hyperref},
\texttt{microtype}, \texttt{booktabs}, \texttt{amsmath} and
\texttt{amsthm}.

The entry point is \texttt{thesis.tex}, and the files are laid out as
follows:
\begin{verbatim}
  thesis.cls   the thesis class (structure, typography)
  cover.sty    the cover page
  thesis.tex      entry point
  frontmatter/    cover metadata and the preliminary parts
  chapters/       the main chapters
  backmatter/     the appendices
  references.bib  your bibliography
\end{verbatim}

\section{Document Structure}
\label{sec:structure}

The document is divided in the usual way:
\begin{itemize}
\item \verb|\frontmatter| --- the cover, dedication, declaration, preface,
  acknowledgements and abstract, numbered in roman numerals;
\item \verb|\mainmatter| --- the chapters, numbered in arabic;
\item \verb|\backmatter| --- the bibliography and the appendices.
\end{itemize}
The preliminary parts use the \verb|frontchapter| environment, which
takes the title you give it and starts an unnumbered, table-of-contents
listed page:
\begin{verbatim}
  \begin{frontchapter}[Declaration]
    ... the declaration text ...
  \end{frontchapter}
\end{verbatim}
(The title may also be given in braces, \verb|\begin{frontchapter}{Declaration}|.)
The abstract is a special page with its own \verb|abstract| environment.
Each preliminary part is kept in its own file under \texttt{frontmatter/}
and simply included from \texttt{thesis.tex}.

\section{Modifications}
\label{sec:modifications}

The class is meant to be used as-is. Adjusting margins, typeface sizes,
line spacing, paragraph and list definitions, and the use of the
\verb|\vspace| command to manually adjust the vertical spacing between
elements of your work should be done with care, as these settings often
need to satisfy the requirements of your institution.
